We develop methods for fitting a viral kinetics model given a set of \vlr{(is there a word more general than "saliva" that I can use here?)} samples collected from $\numpeople$ individuals over $\numdays$ days.
These samples can be arranged into an $\numpeople$ by $\numdays$ matrix $\samplematrix$, where an entry $\samplematrixentry_{i,j}$ is the sample collected from participant $i$ on day $j$.
We say that $\samplematrixentry_{i,j} > 0$ if individual $i$ had viral RNA concentration above the limit of detection (LOD) of our assay on day $j$.
Otherwise, $\samplematrixentry_{i,j} = 0$. 
Observing the value of an entry in the matrix amounts to assaying the corresponding sample.

Data of this kind might emerge in a variety of ways.
Samples originally collected for a study focusing on one pathogen can be arranged into this format and used to fit a model of another pathogen.
\vlr{(Need another example of where this data might come from.)}
Of course, researchers might collect samples from a population with the express purpose of employing the assaying approaches proposed here.

The ultimate goal of this work is to develop efficient algorithms to obtaining a fit model of the viral kinetics of a pathogen, starting from the matrix $\samplematrix$ whose values we have not yet observed. 
There are two main types of algorithms we propose towards this end: the first locates infections within the sample matrix and the second fits a viral kinetics model to those infections.

Unless researchers specifically target individuals who they know are likely to become infected within a fixed period of time (an approach which incurs its own high expenses), $\samplematrix$ will be a sparse matrix.
If we were to assay every sample in $\samplematrix$, most of the results will be negative, offering little information to our model.
That is, only a small proportion of the assays perfomed in the complete assaying case are informative.
This suggests that we can locate the infections within $\samplematrix$ using far fewer than the $\numpeople * \numdays$ assays required by the complete testing approach.
But how many tests are actually needed, and how should we distribute them among the samples?

\subsection{Sample Pooling and Interval Sampling}
The problem of infection location amounts to identifying continuous sequences of positive entries within the rows of $\samplematrix$.
We need not record the entire infection from start to finish, but rather identify at least one of the positive values within it.
This can be achieved by the following algorithm.

Suppose you are given the matrix $\samplematrix$ with unobserved entries.
Assume that the duration of each infection $\li$ is drawn from some continuous positive distribution with practical lower bound $\lo$ such that $\prob(\li < \lo) < c$ for some small constant $c$.
\vlr{(how are we doing this in sims since we're working in discrete time?)}

The algorithm proceeds as follows. \vlr{(a diagram might be useful for this)}
Within each row: pool the samples and apply a single assay to the pool \vlr{(gotta show that we'll stay above LO)}. 
For each row $i$ on which the assay returns a negative result, set $\samplematrixentry_{i,j} = 0$ for all $i$.
This row does not contain an infection and can now be removed from consideration.
For each row on which the assay returned a positive result, individually assay entries at an interval of 